

Για την επεξεργασία που έγινε στο τελευταίο notebook,το \texttt{search data close to sensors.ipynb},
το κύριο κομμάτι επεξεργασίας είναι να επαναφέρουμε αρκετές από τις στήλες στους τύπους δεδομένων που είχαν πριν την μετατροπή του σε json αρχείο,
όπως και να μεταμορφώσουμε κάποιες στήλες σε μορφές που είναι απαραίτητη για την μετέπειτα ανάπτυξη των μηχανικών μοντέλων.
Ακόμα,στρογκυλοποιούμε τις ευκλίδειες αποστάσεις μεταξύ πηγής και των συσκευών αισθητήρων 

Εδώ είναι που γίνεται εξαμόλυνση του θορύβου των \texttt{VOC original} δεδομένων με το Savitzky-Golay φίλτρο.
Αποτελεί συνάρτηση της βιβλιοθήκης scipy σχετική με τα σήματα (scipy.signal).

Σε αντίθεση με τα rolling windows,που διαμορφώνει κάθε σημείο με βάση των μέσο όρο των γειτονικών δεδομένων,
το Savitzky\-Golay φίλτρο,που η μέθοδος στην scipy ονομάζεται \texttt{savgol\_filter},
κάθε σημείο διαμορφώνεται μετά από την πραγματοποίηση συνέλιξης μεταξύ ενός αριθμού γειτονικών δεδομένων 
και ενός πολυωνύμου βαθμού N.

Επιλέξαμε σαν μέγεθος των γειτονικών δεδομένων,δηλαδή του παραθύρου,να είναι 60 και ο βαθμός πολυωνύμου να είναι το 2.
Έκανε καλύτερη εξαμόλυνση του θορύβου σε σχέση με τα rolling windows,διατηρώντας τα πιο χαρακτηριστικά
τοπικά ελάχιστα και τοπικά μέγιστα των σημάτων του κάθε πειράματος.

Το μόνο κακό αυτήν της μετατροπής δεδομένων,πουδεν αλλάζει την γενική μορφή των δεδομένων,είναι ότι σε σημεία
πολύ ξαφνικών μεγάλων αλλαγών,οι τιμές που βρίσκονται στα τοπικό ελάχιστα κατεβαίνουν σε αρνητικούς αριθμούς,
κάτι που έχουμε καταστήσει ότι κανονικά δεν πρέπει να γίνεται.
Επείδη τα σημεία αυτά είναι μικρά σε εύρος,σπάνια, δεν διαφθήρουν πολύ την κανονική μορφή των δεδομένων,καθώς συμφωνούν με τα υπολοιπα δεδομένα και δεν χρησιμοποιούμε
κάποια μέθοδο ή αλγόριθμο που χρειάζεται καθαρά θετικές τιμές,δεν μας ενοχλεί τόσο.

Ακόμα,υπολογίσαμε για κάθε πείραμα,ποια θα ήταν η ιδεατή στιγμή αυξήσης των \texttt{VOC original} για κάθε αισθητήρα,
με βάση την ευκλίδια απόσταση του αισθητήρα από την πηγή ρύπου.
Αυτό σχετίζεται με μία προσπάθεια προσέγγισης του ιδεατού μοντέλο με το οποίο θεωρήσαμε ότι θα δουλεύψουν
τα δεδομένα και οι τιμές των αισθητήρων.Θα το αναλύσουμε πιο μετά όταν μιλήσουμε για τα αποτελέσματα των μηχανικών μοντέλων.


\subsection{Η άντληση πληροφορίων από την οπτικοποίηση}

Θα κάνουμε κάποια σχόλια πρώτα για το πως αναπτύχθηκε η οπτικοποίηση των δεδομένων και έπειτα
θα κάνουμε σχόλια για αποτελέσματα που φάνηκαν.
Επείδη ο αριθμός των σχεδιαγραμμάτων των δεδομένων είναι αρκετά μεγάλος,δεν θα τις προσθέσουμε σε αυτό το κεφάλαιο.

Η βασική συνάρτηση που βρίσκεται πίσω από την δημιουργία σχεδιαγραμμάτων των 
\texttt{timeSeries} του κάθε πειράματος,
είναι η \texttt{printDataBasedOnDate}.

Σε αντίθεση με ότι υπονοεί το όνομα της συνάρτησης για σχεδιασμό μέσω ημερομηνίας,
Η συνάρτηση αυτή σχεδιάζει τα \texttt{timeSeries} δεδομένα των πειράματων,ομαδοποιημένα με βάση
τις διακριτές τιμές οποιαδήποτε στήλη της \texttt{userInputData} επιθυμούμε,δίνοντας το αντίστοιχο όνομα της στήλης.

Όταν μιλάμε ομαδοποίηση με διακριτές τιμές,εννοούμε ότι βάζουμε σε μία κοινή ομάδα όσα 
πειράματα έχουν την ίδια τιμή για το πεδίο,για παράδειγμα όλα τα πειράματα που γίνουν στις
14 Οκτώμβρη,και το κάνουμε αυτό για όλες τις τιμές της στήλης που μας δώθηκε που είναι ξεχωριστές.
Χρησιμοποιούμε τον όρο διακριτές και όχι ξεχωριστές,επειδή δεν είναι ότι έχουν κάποια ξεχωριστή 
ιδιότητα,αλλά ότι είναι διαφορετικές μεταξύ τους.

Αν τις δώσουμε πολλαπλά ονόματα στηλών της \texttt{userInputData},
θα ομαδοποιηθούν με σειρά αντίθετης της σειράς που τα εισάγαμε.
Δηλαδή,θα ξεκινήσουμε την επανάληψη ομαδοποιήσεων αλλάζοντας τις διακριτές τιμές τις τελευταίας στήλης,
δηλαδή της στήλης που δωθήκε πιο δεξιό στην λίστα ονομάτων που δώσαμε στις παραμέτρους της συνάρτησης,
κρατόντας τις άλλες στήλες σταθερές και βλέποντας ποια πειράματα υπάρχουν για αυτές τις διακριτές τιμές
των στηλών της \texttt{userInputData}
Αφού εξαντληθούν όλες οι διακριτές τιμές της τελευραίας στήλης,
θα αλλάξουμε την διακριτή τιμή της προτελευταίας στήλης,το όνομα που είναι δεύτερο πιο δεξία, και θα ξαναεπαναλάβουμε τις διακρητές τιμές της τελευταίας στήλης,
και πάει λέγοντας μέχρι να τελειώσουν οι διακριτές τιμές της προτελευταίας στήλης,οπού μετά πάμε στην αμέσως επόμενη διακριτή τιμή.
Όλο αυτό συνεχίζει μέχρι να δούμε όλες τις δυνατές ομαδοποιήσεις των στηλών που μας δώθηκαν.

Ο λόγος που εστιάζουμε τόσο στην ομαδοποίηση με διακριτές τιμές είναι διότι αποτελεί ένα πολύ
δυνατό εργαλείο διαχήρισης των dataframes της pandas, 
την μέθοδο \texttt{groupby()},το οποίο χρησιμοποιούμε πολύ από εδώ και πέρα ,
για να διαχειριστούμε  κυρίως τα \texttt{timeSeries},αλλά και τα \texttt{userInputData}.

Κρατάμε σταθερά τα χαρακτηριστικά χρώματα που είχαμε και στην grafana,δηλαδή μπλε για το αισθητήρα
με id=0,πράσινο για id=1 και κίτρινο για id=2